\subsection{Clifford group.}

In this section, we assume that A is a field, that E has finite dimension m, and that Q is nondegenerate. We identify E with its canonical image in C(Q).

DEFINITION 2. --- We call the Clifford group of Q (resp. the special Clifford group of Q) the multiplicative group of invertible elements s of C(Q) (resp. C \(^{+}\) (Q)) such that sEs \(^{-1}\) = E.

In this section, we denote by G and G⁺ the Clifford group and the special Clifford group of Q. It is clear that G⁺ = G ∩ C⁺(Q).

THEOREM 4. --- Set, for \(s \in G\) and \(x \in E\) , \(\varphi(s). x = sxs^{-1}\) .

\begin{enumerate}
\def\labelenumi{\alph{enumi})}
\item
  The map \(\varphi\) is a homomorphism from G into the orthogonal group \(\mathbf{O}(\mathrm{Q})\) of Q and its kernel is the set of invertible elements of the center Z of C(Q).
\item
  The set \(E \cap G\) is the set of nonsingular vectors of E; for \(x \in E \cap G, -\varphi(x)\) is the symmetry with respect to the hyperplane orthogonal to x.
\item
  If dim(E) is even, we have \(\varphi(G) = \mathbf{O}(Q)\) , \(\varphi(G^{+})\) has index 2 in \(\mathbf{O}(Q)\) if \(E \neq \{0\}\) , and is equal to \(\mathbf{SO}(Q)\) if A has characteristic \(\neq 2\) .
\item
  If dim(E) is odd (which implies that A has characteristic ≠ 2), we have φ(G) = φ(G⁺) = SO(Q).
\end{enumerate}

Indeed, we have \(\mathrm{Q}(sxs^{-1}) = (sxs^{-1})^2 = sx^2s^{-1} = \mathrm{Q}(x)\) for \(s \in G\) and \(x \in E\) , which shows that \(\varphi(s) \in \mathbf{O}(Q)\) . For \(\varphi(s) = 1\) , it is necessary and sufficient that s commute with the elements of E, that is, belong to the center Z of C(Q). This proves a).

For an element \(x\) of E to belong to G, it must be invertible, that is, it must be a non-singular vector (since \(x^2 = Q(x)\) ). If this is so, we have \(x^{-1} = Q(x)^{-1}x\) , whence, for every \(y \in E\) ,

\[
x y x ^ {- 1} = \mathrm{Q} (x) ^ {- 1} x y x = \mathrm{Q} (x) ^ {- 1} x (\Phi (x, y) - x y) = - (y - \Phi (x, y) \mathrm{Q} (x) ^ {- 1} x)
\]

this proves b) (§ 6, no. 4).

Lemma 5. --- Every element s of G is of the form \(zs'\) , where z is an invertible element of Z and \(s'\) belongs to G ∩ C \(^{+}\) (Q) or to G ∩ C \(^{-}\) (Q); the subgroup G \(^{+}\) has index 2 in G when E ≠ \{0\}.

The second assertion follows evidently from the first, since the nonsingular vectors belong to \(\mathrm{G}\cap\mathrm{C}^{-}(\mathrm{Q})\) . Suppose first that dim(E) is even, and let \(s=t'+t''\) , with \(t'\in\mathrm{C}^{+}(\mathrm{Q})\) and \(t''\in\mathrm{C}^{-}(\mathrm{Q})\) ; by definition we have \(sx=(\varphi(s).x)s\) for all \(x\in E\) ; since \(t'x\) and \((\varphi(s).x)t'\) (resp. \(t''x\) and \((\varphi(s).x)t'')\) are odd (resp. even) elements, we have \(t'x=(\varphi(s).x)t'\) , whence \(s^{-1}t'x=xs^{-1}t'\) for all \(x\in E\) . From this we conclude that \(s^{-1}t'\in Z\) and since dim(E) is even, \(Z=A\) (no. 4, cor. of th. 2) , hence \(t'=as\) , where \(a\in A\) . If \(a\neq0\) , we therefore have \(s=a^{-1}t'\) and \(t'\in G\cap C^{+}(Q)\) ; if a=0, \(s=t''\in G\cap C^{-}(Q)\) and the lemma is proved in this case. If dim(E) is odd, A has characteristic \(\neq2\) , so for every \(s\in G\) , \(\varphi(s)\) is a product of symmetries with respect to nonsingular vectors \(x_{i} (i = 1,2,\ldots,h)\) ( \(\S 6, n^{o} 4\) , prop. 5); if we set \(s' = x_{1}x_{2}\ldots x_{h}\) , one has \(\varphi(s) = \varphi(s')\) , hence \(s = zs'\) , where \(z \in \mathbf{Z}\) , and \(s'\) belongs to \(C^{+}(Q)\) or to \(C^{-}(Q)\) according as \(h\) is even or odd.

Suppose dim(E) is even. Since every element \(u\) of \(\mathbf{O}(Q)\) extends in one and only one way to an automorphism \(\overline{u}\) of C(Q) (prop. 1), and since C(Q) is central simple (th. 2), \(\overline{u}\) is an inner automorphism (chap.~VIII, § 10, no. 1, th. 1). There therefore exists an element \(s\) of G such that \(\varphi(s) = u\) On the other hand, the center of C(Q) is contained in C \(^{+}\) , which implies that \(\varphi(G)/\varphi(G^{+})\) is isomorphic to G/G \(^{+}\) , hence that \(\varphi(G^{+})\) has index 2 in \(\varphi(G) = \mathbf{O}(Q)\) if E ≠ \{0\}. This proves the first two assertions of \(c\) ).

Finally, suppose that A has characteristic ≠ 2. Then every element u of O(Q) is a product of symmetries with respect to hyperplanes orthogonal to nonsingular vectors \(x_{i} (i = 1, \ldots, h) (\S 6, n^{o} 4, \text{prop. } 5)\) ; we therefore have \(u = (-1)^{h} \varphi(x_{1} \ldots x_{h})\) and det \((u) = (-1)^{h}\) . For u to belong to SO(Q), it is necessary and sufficient that h be even, which shows that \(\varphi(G^{+}) \supset \mathbf{SO}(Q)\) . Since SO(Q) has index 2 in O(Q) when \(E \neq \{0\}\) , one has \(\varphi(G^{+}) = \mathbf{SO}(Q)\) If E is of even dimension, which completes the proof of c). On the other hand, if the dimension of E is odd, \(\varphi(G)\) does not contain the orthogonal transformation \(x \to -x\) ; indeed, the latter extends to the principal automorphism \(\alpha\) of C(Q) (no. 1), and \(\alpha\) is not an inner automorphism since the center Z of C(Q) contains a nonzero element of C⁻(Q) (th. 3). We therefore have \(\varphi(G) \neq \mathbf{O}(Q)\) , and, since \(\varphi(G) \supset \varphi(G^{+}) \supset \mathbf{SO}(Q)\) and since SO(Q) is of index 2 in O(Q), we have \(\varphi(G) = \varphi(G^{+}) = \mathbf{SO}(Q)\) This proves d). QED.

The subgroup \(\varphi(G^{+})\) of \(\mathbf{O}(Q)\) , which is of index 2 if \(E \neq \{0\}\) , is called the group of rotations of E, and its elements are called rotations; it is denoted by \(\mathbf{O}^{+}(Q)\) Note that, if A is not of characteristic 2, we have \(\mathbf{O}^{+}(Q) = \mathbf{S}\mathbf{O}(Q)\) (cf.~exerc. 9).

PROPOSITION 4. — Let β be the principal antiautomorphism of C(Q) (no. 1). For every \(s \in G^{+}\) , \(\beta(s)s\) is a scalar. The map N : \(s \to \beta(s)s\) is a homomorphism of \(G^{+}\) in the multiplicative group A* of nonzero elements of A.

Indeed, for \(s \in G^{+}\) , one has \(sEs^{-1} = E\) , whence \(\beta(s)^{-1}E\beta(s) = E\) , which shows that \(\beta(s) \in G^{+}\) . Since \(sx = (\varphi(s).x)s\) for all \(x \in E\) , one has \(x\beta(s) = \beta(sx) = \beta(s)(\varphi(s).x)\) and consequently \(\beta(s)sx = \beta(s)(\varphi(s).x)s = x\beta(s)s\) , which implies that \(\beta(s)s\) belongs to the center of C(Q). Since, moreover, \(\beta(s)s\) belongs to C+(Q), \(\beta(s)s\) is a scalar (Ths. 2 and 3). Finally we have \(\beta(st)st = \beta(t)\beta(s)st = \beta(s)s\beta(t)t\) , that is, N(st) = N(s)N(t) for s, t in G+. QED.

The scalar \(N(s) = \beta(s)s (s \in G^{+})\) is called the spinorial norm of s. One denotes by \(G_{0}^{+}\) and one calls the reduced Clifford group the kernel of the homomorphism N. The image \(\varphi(G_{0}^{+})\) is denoted \(\mathbf{O}_{0}^{+}(Q)\) and is called the reduced orthogonal group of Q. Since the kernel of the restriction of \(\varphi\) to \(G^{+}\) is the set of even invertible elements of the center of C(Q) (Th. 4) and hence is identified with A* (Ths. 2 and 3), \(\varphi(G^{+})/\mathbf{O}_{0}^{+}(Q)\) is isomorphic to \(G^{+}/A^{*}G_{0}^{+}\) , hence also with \(N(G^{+})/N(A^{*})\) , and in particular is commutative. It is clear that \(N(A^{*})\) is the subgroup \((A^{*})^{2}\) of squares of elements of A*. If Q has index \textgreater0, there exist, for every \(a \in A^{*}\) , two elements x and y of E such that \(Q(x) = a\) and \(Q(y) = 1 (\S 4, n^{\circ} 2, \text{prop. } 4)\) ; since \(xy \in G^{+}\) and \(N(xy) = Q(x)Q(y) = a\) ; this shows that \(N(G^{+}) = A^{*}\) , hence that \(\varphi(G^{+})/\mathbf{O}_{0}^{+}(Q)\) is isomorphic to \(A^{*}/(A^{*})^{2}\) .

Exercises. --- ¶ 1) Prove Corollaries 3 and 4 of Theorem 1 of No. 3 when E₁ and E₂ are any two complementary submodules of E. (First establish Cor. 4: to do this, show that the tensor product C(Q₁) ⊗ C(Q₂) is endowed with an algebra structure by the sign convention made in the statement of Cor. 4, and that this algebra S solves the same universal mapping problem as C(Q); for this, consider, for every linear map f of E into an algebra D over A such that (f(x))² = Q(x)·1, the homomorphism \(\bar{f}_i\) of C(Qᵢ) into D such that \(\bar{f}_i(1) = 1\) , \(\bar{f}_i(\rho_{Q_i}(x_i)) = f(x_i)\) for \(x_i \in \mathrm{E}_i\) ( \(i = 1, 2\) ) and prove that there exists a homomorphism \(\bar{f}\) of the algebra S into D such that

\[
\bar {f} (z _ {1} \otimes z _ {2}) = \bar {f} _ {1} (z _ {1}) \bar {f} _ {2} (z _ {2})
\]

for \(z_{i}\in\mathrm{C}(\mathrm{Q}_{i})\) ( \(i=1,2\) ). To then prove Cor. 3, consider the quadratic form Q' which is the external direct sum of Q \(_{1}\) and Q \(_{2}\) (§ 3, No. 4), and note that one has Q'(x)=Q(x)+F(x,x), F being the bilinear form defined by F(x \(_{1}\) +x \(_{2}\) , y \(_{1}\) +y \(_{2}\) ) = -Φ(x \(_{1}\) , y \(_{2}\) ) for x \(_{i}\) , y \(_{i}\) in E \(_{i}\) ( \(i=1,2\) ); then use prop. 3.)

¶ 2) Suppose that E is the direct sum of two orthogonal submodules E \(_{1}\) , E \(_{2}\) and denote by Q \(_{i}\) the restriction of Q to E \(_{i}\) (i = 1, 2);

suppose moreover that there exists \(u \in \mathrm{C}^{+}(\mathrm{Q}_{2})\) such that \(u^{2} = a.1\) , \(a\) being invertible in A, and that \(u\rho_{\mathrm{Q}_{2}}(x_{2}) = -\rho_{\mathrm{Q}_{2}}(x_{2})\) for all \(x_{2} \in \mathrm{E}_{2}\) . Show that there exists an isomorphism \(\varphi\) of C(Q) onto the tensor product \(\mathrm{C}(a\mathrm{Q}_{1}) \otimes \mathrm{C}(\mathrm{Q}_{2})\) (as defined in Chap.~III, § 3, no. 1), having the property that for every \(x = x_{1} + x_{2}\) ( \(x_{i} \in \mathrm{E}_{i}, i = 1, 2\) ), one has

\[
\varphi (\rho_ {q} (x)) = \rho_ {a q _ {1}} (x _ {1}) \otimes u ^ {- 1} + 1 \otimes \rho_ {q _ {2}} (x _ {2}).
\]

(Prove the existence of the homomorphism \(\varphi\) as a consequence of the universal property of C(Q). On the other hand, there is a homomorphism \(g_{1}\) of C(aQ \(_{1}\) ) into C(Q) such that \(g_{1}(\rho_{aQ_{1}}(x_{1})) = h_{2}(u)\rho_{Q_{1}}(x_{1})\) , denoting by \(h_{2}\) the canonical homomorphism of C(Q \(_{2}\) ) into C(Q); then observe that \(g_{1}(z_{1})\) and \(h_{2}(z_{2})\) commute for \(z_{1} \in\) C(aQ \(_{1}\) ) and \(z_{2} \in\) C(Q \(_{2}\) ), and deduce from this the existence of an inverse homomorphism of \(\varphi\) .)

In which case does this result apply when A is a field of characteristic ≠ 2 (use Remark 2 following the cor. of th. 2)?

\begin{enumerate}
\def\labelenumi{\arabic{enumi})}
\setcounter{enumi}{2}
\tightlist
\item
  \begin{enumerate}
  \def\labelenumii{\alph{enumii})}
  \tightlist
  \item
    With the notations of no. 2, let \(x_{i}\) ( \(1 \leqslant i \leqslant n\) ) be elements of E such that \(f(x_{i}) = 0\) ; show that one has \(i_{j}(x_{1} \otimes x_{2} \otimes \cdots \otimes x_{n}) = 0\) . In particular, if F is a bilinear form on E such that F( \(x_{i}, x_{j}\) ) = 0 for \(i > j\) , one has \(i_{x_{1}}^{\mathrm{F}}(x_{2} \otimes x_{3} \otimes \cdots \otimes x_{n}) = 0\) .
  \end{enumerate}
\end{enumerate}

\begin{enumerate}
\def\labelenumi{\alph{enumi})}
\setcounter{enumi}{1}
\tightlist
\item
  With the notations of Prop. 3 of no. 3, let \(x_i\) ( \(1 \leqslant i \leqslant n\) ) be elements of E such that \(F(x_i, x_j) = 0\) for \(i > j\) . Show that we have \(\overline{\lambda}_{\mathrm{F}}(\rho_{\mathrm{Q}'}(x_1) \ldots \rho_{\mathrm{Q}'}(x_n)) = \rho_{\mathrm{Q}}(x_1) \ldots \rho_{\mathrm{Q}}(x_n)\) (using \(a\) ) and formula (10) of no. 2).
\item
  We assume that A is a field of characteristic \(\neq 2\) ; for every quadratic form Q on E, let \(\mu_{\mathrm{Q}}\) be the map \(\overline{\lambda}_{\mathrm{F}}\) of C(Q) onto \(\wedge\) E corresponding to \(F(x, y) = \frac{1}{2}\Phi(x, y)\) . Show that if the vectors \(x_i\) ( \(1 \leqslant i \leqslant n\) ) are pairwise orthogonal, we have
\end{enumerate}

\[
\mu_ {\mathrm{Q}} (x _ {1} x _ {2} \dots x _ {n}) = x _ {1} \wedge x _ {2} \wedge \dots \wedge x _ {n}.
\]

Deduce that if \(y_{i}\) ( \(1 \leqslant i \leqslant n\) ) are \(n\) arbitrary vectors, we have

\[
n! \mu_ {0} ^ {- 1} (y _ {1} \wedge y _ {2} \wedge \dots \wedge y _ {n}) = \sum_ {\sigma \in \mathfrak {S} _ {n}} \varepsilon_ {\sigma} y _ {\sigma (1)} y _ {\sigma (2)} \dots y _ {\sigma (n)}.
\]

\begin{enumerate}
\def\labelenumi{\arabic{enumi})}
\setcounter{enumi}{3}
\tightlist
\item
  Let \(C_{h}\) be the submodule of C(Q) generated by the products of k elements of \(\rho_{Q}(E)\) for \(0 \leqslant k \leqslant h\) . Show that the map
\end{enumerate}

\[
(x _ {1}, \dots , x _ {h}) \rightarrow \rho_ {\mathrm{Q}} (x _ {1}) \dots \rho_ {\mathrm{Q}} (x _ {h})
\]

defines by passing to quotients an alternating multilinear map from \(E^{h}\) into \(C_{h}/C_{h-1}\) . From this we deduce a linear map \(\pi_{h}\) of \(\bigwedge E\) into \(C_{h}/C_{h-1}\) ; show that \(\pi_{h}\) is an isomorphism when E is a free module. If f is an orthogonal transformation, we have \(C(f)\circ\pi_{h}=\pi_{h}\circ(\bigwedge f)\) . When A is a field of characteristic \(\neq2\) , show that \(\pi_{h}\) is the restriction of the map \(\mu_{Q}^{-1}\) defined in Exerc. 3 c).

\begin{enumerate}
\def\labelenumi{\arabic{enumi})}
\setcounter{enumi}{4}
\tightlist
\item
  With the notation of no. 2, consider the map \(i_{f}\) of \(\wedge\) from E into itself. Show that one has
\end{enumerate}

\[
i _ {f} \left(x _ {1} \wedge \dots \wedge x _ {h}\right) = \sum_ {i = 1} ^ {h} (- 1) ^ {i - 1} f \left(x _ {i}\right) \left(x _ {1} \wedge \dots \wedge x _ {i - 1} \wedge x _ {i + 1} \wedge \dots \wedge x _ {h}\right)
\]

and deduce that \(i_{f}\) is none other than the right interior product \(z \to z \sqcup f\) (chap.~III, § 8, no. 4, def. 2).

\begin{enumerate}
\def\labelenumi{\arabic{enumi})}
\setcounter{enumi}{5}
\item
  Show that, if E admits an orthogonal basis \((e_{1}, e_{2})\) for Q, and if we set \(\alpha_{i} = \mathrm{Q}(e_{i})\) For (i = 1, 2), the algebra C(Q) is isomorphic to the algebra of quaternions over A corresponding to the pair \((\alpha_{1}, \alpha_{2})\) .
\item
  Let A be a maximal ordered field, E a vector space of even dimension 2r over A, Q a nondegenerate quadratic form on E, positive or negative. Show that, if Q is positive, C(Q) is isomorphic to a matrix algebra over A when \(r(r-1)/2\) is even, to a matrix algebra over the quaternion algebra over A when \(r(r-1)/2\) is odd. If Q is negative, C(Q) is isomorphic to a matrix algebra over A when \(r(r+1)/2\) is even, to a matrix algebra over the quaternion algebra over A when \(r(r+1)/2\) is odd (use Exercises 2 and 6). What is the structure of C⁺(Q) in these different cases?
\item
  Let A be the ring \(\mathbf{Z}/(4)\) a left A-module \(\mathbf{Z}/(2)\) ; if we set \(Q(0)=0\) , \(Q(u)=1\) for the unique element \(u\neq0\) of E, Q is a quadratic form on E. Show that the A-modules C(Q) and \(\wedge\) They are not isomorphic (one will prove that C(Q) is isomorphic to the direct sum of two modules isomorphic to E).
\end{enumerate}

¶ 9) Suppose A is a field of characteristic 2, E a vector space of finite even dimension 2r, and Q a nondegenerate quadratic form on E.

\begin{enumerate}
\def\labelenumi{\alph{enumi})}
\tightlist
\item
  Let \((e_i)\) a symplectic basis (§ 5, no. 1) of E for the alternating bilinear form \(\Phi\) associated with Q. Show that the element
\end{enumerate}

\[
z = e _ {1} e _ {2} + e _ {3} e _ {4} + \dots + e _ {2 r - 1} e _ {2 r}
\]

of C \(^{+}\) (Q), together with the identity element, forms a basis of the center Z of C \(^{+}\) For Z to be a direct sum of two fields, it is necessary and sufficient that the element

\[
\Delta (Q) = Q \left(e _ {1}\right) Q \left(e _ {2}\right) + Q \left(e _ {3}\right) Q \left(e _ {4}\right) + \dots + Q \left(e _ {2 r - 1}\right) Q \left(e _ {2 r}\right)
\]

(called the pseudo-discriminant of Q with respect to the symplectic basis \((e_{i})\) ) is of the form \(\lambda^{2} + \lambda\) , with \(\lambda \in A\) .

\begin{enumerate}
\def\labelenumi{\alph{enumi})}
\setcounter{enumi}{1}
\tightlist
\item
  Let \(u\) a similarity for \(\Phi\) ( \(\S 6\) , no. 5) with multiplier \(\mu(u)\) , and let \(Q_{1}(x) = Q(u(x))\) . We set
\end{enumerate}

\[
u (e _ {2 i - 1}) = \sum_ {j = 1} ^ {r} a _ {i j} e _ {2 j - 1} + \sum_ {j = 1} ^ {r} b _ {i j} e _ {2 j},
\]

\[
u (e _ {2 i}) = \sum_ {j = 1} ^ {r} c _ {i j} e _ {2 j - 1} + \sum_ {j = 1} ^ {r} d _ {i j} e _ {2 j},
\]

and \(\mathrm{Q}(e_{2i - 1}) = \alpha_i, \mathrm{Q}(e_{2i}) = \beta_i (1 \leqslant i \leqslant r)\) . Show that we have

\[
\Delta (\mathrm{Q} _ {1}) = (\mu (u)) ^ {2} \Delta (\mathrm{Q}) + (\mathrm{D} (u)) ^ {2} + \mu (u) \mathrm{D} (u)
\]

where

\[
\mathbf {D} (u) = \sum_ {i, j} (\alpha_ {j} a _ {i j} c _ {i j} + \beta_ {j} b _ {i j} d _ {i j} + b _ {i j} c _ {i j})
\]

(Dickson invariant of u with respect to the basis \((e_{i})\) ). (Note that the element

\[
(\mu (u)) ^ {- 1} (u (e _ {1}) u (e _ {2}) + u (e _ {3}) u (e _ {4}) + \dots + u (e _ {2 r - 1}) u (e _ {2 r}))
\]

belongs to Z.) For \(u\) to be a similarity for Q (§ 4, exerc. 9) with multiplier \(\mu(u)\) , it is necessary and sufficient that D(u) = 0 or D(u) = \(\mu(u)\) ; the similarities for Q such that D(u) = 0 are called direct.

\begin{enumerate}
\def\labelenumi{\alph{enumi})}
\setcounter{enumi}{2}
\tightlist
\item
  Show that, if v is a similarity for \(\Phi\) , u a similarity for Q, one has
\end{enumerate}

\[
\mathrm{D} (u v) = \mu (v) \mathrm{D} (u) + \mu (u) \mathrm{D} (v)
\]

(consider the Dickson invariant of \(\varrho\) with respect to the symplectic basis formed by the \(u(e_{2i-1})\) and \((\mu(u))^{-1}u(e_{2i})\) for \(1 \leqslant i \leqslant r\) ). Deduce from this that the direct similarities for Q form a distinguished subgroup of index 2 in the group of similarities for Q.

\begin{enumerate}
\def\labelenumi{\alph{enumi})}
\setcounter{enumi}{3}
\item
  If \(u\) is the symmetry with respect to the hyperplane orthogonal to a nonsingular vector in E, show that \(\mathrm{D}(u)=1\) . Deduce from this that the group \(\varphi(\mathrm{G}^{+})\) (notations of no. 5) is the group \(\mathbf{SO}(\mathrm{Q})\) defined in exerc. 28 c) of § 6.
\item
  Let \(u \in \mathbf{O}(Q)\) , and suppose \(\Lambda\) algebraically closed. Show that, for \(u \in \mathbf{SO}(Q)\) , it is necessary and sufficient that the number of elementary divisors of the module \(E_u\) be even. (With the notations of exerc. 15 of § 5, first note that if \(p, q\) are two distinct irreducible factors of the minimal polynomial of \(u\) , the number of indecomposable submodules of which \(G(p, q)\) is the direct sum is equal to the number of indecomposable submodules of which \(G(q, p)\) is the direct sum. Note on the other hand that \(G(p, p) = \{0\}\) except for \(p = X - 1\) . Finally prove that one may restrict to the case where \(E_u\) is equal to \(G(p, p)\) (with \(p = X - 1\) ) and is indecomposable. There is then a symplectic basis ( \(e_i\) ) of E such that \(e_1, e_3, \ldots, e_{2k-1}\) form a basis of \((p(u))^{2r-k}(E)\) for \(1 \leqslant k \leqslant r\) ; show that \(e_1, e_3, \ldots, e_{2r-3}\) are singular vectors, and conclude that \(D(u) = 1\) .)
\end{enumerate}

¶ 10) Suppose A is a field, E a finite-dimensional vector space, Q a degenerate quadratic form on E; let M be a subspace complementary to E \(^{0}\) in E, let B be the (semisimple) Clifford algebra of the restriction of Q to M.

\begin{enumerate}
\def\labelenumi{\alph{enumi})}
\item
  We first assume that A has characteristic ≠ 2. Let L be the Clifford algebra of the restriction of Q to E⁰ (isomorphic to ∧ E⁰), and R₀ its radical (the ideal generated in L by E⁰, of codimension 1 in L); show that the radical R of C(Q) is obtained (up to isomorphism) by defining the algebra structure on \(B \otimes_{A} R_{0}\) as in cor. 4 of th. 1, that C(Q)/R is isomorphic to B, and that C(Q) is the direct sum of B and R.
\item
  Assume A has characteristic 2. Let F be the subspace of \(\mathbf{E}^0\) consisting of singular vectors \(x \in \mathbf{E}^0\) and let N be a complement of F with respect to \(\mathbf{E}^0\) . If \((a_i)_{1 \leqslant i \leqslant d}\) is a basis of N, and \(\mathrm{Q}(a_i) = \alpha_i\) the elements \(\alpha_i^{1/2}\) in an algebraic closure of A, are linearly independent over A. Let \((\alpha_i^{1/2})_{1 \leqslant i \leqslant e}\) be a 2-basis of the field \(\mathrm{A}_1 = \mathrm{A}(\alpha_1^{1/2}, \ldots, \alpha_d^{1/2})\) on A (Chap.~V, § 8, exerc. 1), and let \(h = \dim F\) . If \(\mathrm{B}_1\) is the central simple algebra \(\mathrm{B} \otimes_{A} \mathrm{A}_1\) , show that \(\mathrm{C(Q)}\) is isomorphic to the algebra \(\mathrm{B}_1 \otimes_{A_1} \mathrm{L}_1\) , where \(\mathrm{L}_1\) is the exterior algebra of a vector space of dimension \(h + d - e\) on \(\mathrm{A}_1\) . If \(\Re_1\) is the radical of \(\mathrm{L}_1\) (of codimension 1 (over \(\mathrm{A}_1\) ) in \(\mathrm{L}_1\) ), the radical \(\Re\) of \(\mathrm{C(Q)}\) is isomorphic to the algebra \(\mathrm{B}_1 \otimes_{A_1} \Re_1\) , \(\mathrm{C(Q)/R}\) is isomorphic to \(\mathrm{B}_1\) , and \(\mathrm{C(Q)}\) is a direct sum of \(\mathrm{B}_1\) and of \(\Re\) .
\item
  The hypotheses being those of \(b\) ; one assumes moreover that \(h=0\) , \(d=1\) , \(e=0\) (which implies that dim M is even). If Q \(_{0}\) is the restriction of Q to M, show that C \(^{+}\) (Q) is isomorphic to C(Q \(_{0}\) ).
\end{enumerate}

11) a) Under the hypotheses and notation of no. 5, show that N(G⁺) is the subgroup of A* generated by the products Q(x)Q(y), where x and y range over the set of non-isotropic vectors of E. (Reduce to the case A ≠ F₂; use prop. 5 and exerc. 28 c) of § 6, as well as exerc. 9 d) of § 9.) Case where Q has index ≥ 1.

\begin{enumerate}
\def\labelenumi{\alph{enumi})}
\setcounter{enumi}{1}
\item
  We assume in addition that \(\mathbf{A} \neq \mathbf{F}_2\) that E has dimension \(n \geqslant 2\) and that Q has index \(>0\) . Show that \(\mathbf{O}_0^+(\mathbf{Q})\) is the commutator subgroup of the group \(\mathbf{O}(\mathbf{Q})\) . (Reduce to the case \(n = 2\) using Exercises 17 c) and 28 e) of § 6.)
\item
  We keep the hypotheses of \(b\) , and we suppose in addition that A has characteristic \(\neq 2\) and that E is of even dimension. For the automorphism \(x \to -x\) of E belong to \(\mathbf{O}_0^+(\mathbf{Q})\) it is necessary and sufficient that the discriminant of Q (with respect to any basis of E) be a square.
\end{enumerate}

Let a be an invertible element of A; show that there exists one and only one isomorphism \(\theta_{a}\) of C \(^{+}\) (Q) onto C \(^{+}\) (aQ) such that

\[
\theta_ {a} (\rho (x) \rho (y)) = a ^ {- 1} \rho_ {1} (x) \rho_ {1} (y)
\]

for all \(x, y\) in E ( \(\rho\) and \(\rho_{1}\) denoting the canonical maps from E into C(Q) and C(aQ), respectively). Deduce that if A is a field of characteristic \(\neq 2\) , E a vector space of odd dimension and Q a non-degenerate quadratic form, C(Q) and C(aQ) are isomorphic (cf.~exerc. 7).

\begin{enumerate}
\def\labelenumi{\alph{enumi})}
\setcounter{enumi}{1}
\tightlist
\item
  Assume that A is a field, E a vector space of even dimension 2r \textgreater{} 0, Q a nondegenerate quadratic form. Let u be a similarity relative to Q; show that there exists one and only one A-automorphism \(\overline{u}\) of the algebra \(\mathrm{C}^{+}(\mathrm{Q})\) such that for \(1 \leqslant h \leqslant r\) and \(x_{i} \in E (1 \leqslant i \leqslant 2h)\) one has
\end{enumerate}

\[
\overline {{{u}}} (x _ {1} x _ {2} \dots x _ {2 h}) = \mu^ {- h} u (x _ {1}) u (x _ {2}) \dots u (x _ {2 h})
\]

where \(\mu\) denotes the multiplier of \(u\) . For \(\bar{u}\) to be an inner automorphism, it is necessary and sufficient that \(u\) be a direct similarity (§ 6, no. 5 and § 9, exerc. 9 b)). One assumes in addition \(r \geqslant 2\) ; then, for \(\bar{u}\) to be the identity, it is necessary and sufficient that \(u\) be a homothety.

Show that the automorphism \(\bar{u}\) of C \(^{+}\) (Q) is the restriction to C \(^{+}\) (Q) of an inner automorphism of C(Q).

¶ 13) Suppose that A is a field of characteristic ≠ 2, E a vector space of even dimension 2r. \textgreater{} 0, Q a nondegenerate quadratic form. We denote by \(E_{h}\) the inverse image of \(\Lambda E\) under the isomorphism \(\mu_{Q}\) of exerc. 3 c), so that \(E_{h}\) is the subspace of C(Q) generated by the products \(x_{1}x_{2}\ldots x_{h}\) , where the \(x_{i}\) ( \(1 \leqslant i \leqslant h\) ) are pairwise orthogonal.

\begin{enumerate}
\def\labelenumi{\alph{enumi})}
\tightlist
\item
  Let \(\alpha\) an element of A, \(s\) an element of E \(_2\) . Show that, if \((\alpha + s)^2 \in \mathbf{A}\) , then either \(s = 0\) , or else \(\alpha = 0\) and \(s = xy\) , where \(x\) and \(y\) are two orthogonal vectors. (Note that if \(t = \mu_Q(\alpha + s)\) , \(t \wedge t\)
\end{enumerate}

necessarily belongs to A + ∧ E, by expressing s using an orthogonal basis of E; deduce that one necessarily has t = β + x ∧ y with β ∈ A, x, y in E, using § 5, no. 1, cor. 2 of th. 1).

\begin{enumerate}
\def\labelenumi{\alph{enumi})}
\setcounter{enumi}{1}
\item
  Let \((x, y)\) , \((u, v)\) two pairs of orthogonal non-isotropic vectors in E, \(\mathrm{P}_{xy}\) and \(\mathrm{P}_{uv}\) the planes \(Ax + Ay\) , \(Au + Av\) respectively. For one to have \((xy)(uv) = -(uv)(xy)\) in C(Q), it is necessary and sufficient that \(\mathrm{P}_{xy} + \mathrm{P}_{uv}\) be a non-isotropic subspace of dimension 3, in which \(\mathrm{P}_{xy}\) and \(\mathrm{P}_{uv}\) are weakly orthogonal (§ 3, exerc. 11).
\item
  Let g be an A-automorphism of \(\mathrm{C}^{+}(\mathrm{Q})\) sending \(E_{2}\) into itself; show that there exists a similarity u relative to Q such that \(g = \bar{u}\) (exerc. 12 b)). (Let \((e_{i})_{1 \leqslant i \leqslant 2r}\) be an orthogonal basis of E; using a) and b), show that there exists an orthogonal basis \((x_{i})\) of E such that \(g(e_{1}e_{i}) = x_{1}x_{i}\) for \(2 \leqslant i \leqslant 2r\) .)
\end{enumerate}

\begin{enumerate}
\def\labelenumi{\arabic{enumi})}
\setcounter{enumi}{13}
\tightlist
\item
  We assume that A is a field of characteristic \(\neq 2\) E a vector space of dimension \(n\) on A, Q and Q₁ two non-degenerate quadratic forms on E; let \(\tilde{\Delta}\) (resp. \(\tilde{\Delta}_{1}\) ) be the class of the discriminant \(\Delta\) of Q (resp. of the discriminant \(\Delta_{1}\) of Q \(_{1}\) ) with respect to a basis of E, in the group A \(^{*}\) /(A \(^{*}\) )², a class that does not depend on the basis chosen in E. We assume \(n = 2\) .
\end{enumerate}

\begin{enumerate}
\def\labelenumi{\alph{enumi})}
\item
  Show that if \(\tilde{\Delta} = \tilde{\Delta}_{1}\) and if the Clifford algebras C(Q) and C(Q \(_{1}\) ) are isomorphic, then Q and Q \(_{1}\) are equivalent (consider on C(Q) the quadratic form \(z \to z\bar{z}\) , where \(z \to \bar{z}\) is the unique involutive antiautomorphism of C(Q) whose set of invariants is the center of C(Q) (cf.~exerc. 6 and chap.~VIII, § 11, exerc. 5 e)); apply Witt's theorem).
\item
  For C(Q) to be isomorphic to the matrix algebra \(\mathbf{M}_{2}(\mathrm{A})\) it is necessary and sufficient that, for at least one \(x \neq 0\) in E, there exists \(y\) , \(z\) in E such that Q(x) + Q(y)Q(z) = 0; if this is the case, this property holds for every \(x \neq 0\) in E.
\end{enumerate}

¶ 15) We keep the notations and general hypotheses of exerc. 14, but we assume n = 3.

\begin{enumerate}
\def\labelenumi{\alph{enumi})}
\item
  Show that \(\mathrm{C}^{+}(\mathrm{Q})\) is isomorphic to a quaternion algebra over A and that \(\beta\) is the antiautomorphism \(z\to \bar{z}\) of this algebra whose set of invariants is the center of \(\mathrm{C}^{+}(\mathrm{Q})\) ; if P is the subspace of \(\mathrm{C}^{+}(\mathrm{Q})\) consisting of pure quaternions (that is, those such that \(z = -\bar{z}\) ; chap.~VIII, § 11, exerc. 6), the restriction to P of the quadratic form \(z\to z\bar{z}\) is equivalent to \(\lambda Q\) , where \(\lambda \in A\) . Deduce that, in order for \(C^{+}(Q)\) to be a field, it is necessary and sufficient that \(Q\) be of index 0.
\item
  Show that if \(\tilde{\Delta} = \tilde{\Delta}_{1}\) and if the Clifford algebras C(Q) and C(Q \(_{1}\) ) are isomorphic, then Q and Q \(_{1}\) are equivalent. (First consider the case where - \(\Delta\) is a square in K, and show that in this case C \(^{+}\) (Q) and C \(^{+}\) (Q \(_{1}\) ) are isomorphic; then reason as in exerc. 14 a), using a) and exerc. 6 of chap.~VIII, § 11. In the general case, use exerc. 12 a)).
\item
  Show that the special Clifford group G \(^{+}\) (for the form Q) is identical to the group of invertible elements of C \(^{+}\) (Q). (If ( \(e_{1}\) , \(e_{2}\) , \(e_{3}\) ) is an orthogonal basis of E, and if \(j = e_{1}e_{2}e_{3}\) in C(Q), note that \(x \to xj\) is an isomorphism of vector spaces from E onto P).
\item
  Deduce from a) and c) that if Q has index 1, the group of rotations \(\mathbf{O}^{+}(\mathrm{Q})\) is isomorphic to the projective group \(\mathbf{PGL}_{2}(\mathrm{A})\) (chap.~II, \(2^{\mathrm{e}}\) éd., App. III, no. 6).
\end{enumerate}

We keep the general hypotheses and notations of exerc. 14, but we assume n = 4.

\begin{enumerate}
\def\labelenumi{\alph{enumi})}
\item
  Give an example where \(\tilde{\Delta} = \tilde{\Delta}_{1}\) and where \(C(Q)\) and \(C(Q_{1})\) are isomorphic, but where \(Q\) and \(Q_{1}\) are not equivalent (cf.~exerc. 7).
\item
  Let \((e_i)_{1\leqslant i\leqslant 4}\) an orthogonal basis of E for Q, \(Q_0\) the restriction of Q to the hyperplane H = Ae \(_1\) + Ae \(_2\) + Ae \(_3\) Show that, if Z is the center of C \(^+\) (Q), the algebra C \(^+\) (Q) is isomorphic to the tensor product Z ⊗ \(_A\) C \(^+\) (Q \(_0\) ). For every z ∈ C \(^+\) (Q), one has β(z)z ∈ Z ; for z to belong to the special Clifford group G \(^+\) it is necessary and sufficient that z be invertible and that β(z)z ∈ A.1. Deduce from this that the group O \(_0^+\) (Q) is isomorphic to the quotient by \{1,-1\} of the group of elements z ∈ Z ⊗ \(_A\) C \(^+\) (Q \(_0\) ) such that β(z)z = 1.
\item
  We assume that \(\Delta\) is not a square in A (which implies, by Witt's theorem, that the index of Q is 0 or 1). If \(Q_0'\) is the quadratic form obtained from \(Q_0\) by extension to \(A' = A(\sqrt{\Delta})\) of the field of scalars, deduce from \(b\) ) that \(\mathbf{O}_0^+(Q)\) is isomorphic to \(\mathbf{O}_0^+(Q_0')\) . In particular, if Q has index 1, \(\mathbf{O}_0^+(Q)\) is the commutator subgroup of \(\mathbf{O}(Q)\) and is isomorphic to \(\mathbf{PSL}_2(A')\) (cf.~exerc. 15 \(d\) ) and chap.~III, § 7, exerc. 8).
\item
  We assume that \(\Delta\) is a square in A (which implies, by Witt's theorem, that Q has index 0 or 2) and that \(Q(e_4) = 1\) . If one sets \(j = e_1e_2e_3\)  \(x \in E\) can be written in only one way \(x = \alpha e_4 + jz\) , where \(\alpha \in A\) and \(z\) is a pure quaternion (exerc. 15 a)) in \(L = C^+(Q_0)\) ; if we set \(\psi(x) = \alpha + z\) , \(\psi\) is a vector space isomorphism from E onto L. Let \(Z = Ac' + Ac''\) , where \(c'\) and \(c''\) are the two orthogonal idempotents in Z; every invertible element \(s \in C^+(Q)\) is written in a unique way \(s = uc' + vc''\) , where \(u\) and \(\nu\) belong to L; so that \(s\) belongs to the special Clifford group \(G^+\) , it is necessary and sufficient that \(u\bar{u} = v\bar{v}\) , and we then have \(\psi(sxs^{-1}) = u\psi(x)v^{-1}\) for all \(x \in E\) Deduce from this that the quotient of \(\mathbf{O}_0^+(Q)\) by its center (which is a group with 2 elements, cf.~exerc. 11 c)) is isomorphic to the product \(\mathbf{O}_0^+(Q_0) \times \mathbf{O}_0^+(Q_0)\) ; in particular, if Q has index 2, this quotient group is isomorphic to \(\mathbf{PSL}_2(A) \times \mathbf{PSL}_2(A)\) .
\end{enumerate}

\begin{enumerate}
\def\labelenumi{\arabic{enumi})}
\setcounter{enumi}{16}
\tightlist
\item
  Let K be a commutative field of characteristic ≠ 2, A the field \(\mathrm{K}(\mathrm{X}_{n})_{n\in\mathbb{N}}\) of rational fractions over K, with respect to a countable family of indeterminates (chap.~IV, § 3, n° 1). Let E be a vector space over A, having a countable basis \((e_{n})_{n\in\mathbb{N}}\) , and let Φ be a symmetric bilinear form on E, for which \((e_{n})\) is an orthogonal basis and such that \(\Phi(e_{n}, e_{n}) = \mathrm{X}_{n}\) for all \(n \in N\) . If one sets \(\mathrm{Q}(x) = \Phi(x, x)\) , show that the Clifford algebra C(Q) is a field (cf.~chap.~VIII, § 12, exerc. 14).
\end{enumerate}

